\begin{center}
    {\LARGE \textbf{2026无机化学与物理化学}} \\[2ex]
    {\large 时量180分钟 \quad 满分100分}
\end{center}

\vspace{0.5cm}
\noindent\textbf{注意事项:}
\begin{enumerate}[label=\arabic*., parsep=0pt, itemsep=0pt]
    \item 答题前,考生先将自己的姓名、学号、考场号、座位号填写在答题卡上,将条形码准确粘贴在考生信息条形码粘贴区。
    \item 考试结束后,将本试卷与答题纸一并收回。
    \item 回答各个问题时,将答案写在答题纸上,写在本试卷上无效。
\end{enumerate}

\vspace{1cm}
\begin{center}
    {\Large \textbf{无机化学部分(共75分)}}
\end{center}

\vspace{0.5cm}
\noindent\textbf{一、选择题（每题2分，共20分）}

\begin{enumerate}[label=\arabic*., itemsep=1.5ex]
    \item 某基态原子的第六电子层只有2个电子，则该原子的第五层中电子数可以为
    \begin{tasks}(4)
        \task 8个
        \task 8个或18个
        \task 8$\sim$18个
        \task 18$\sim$32个
    \end{tasks}

    \item 下列正盐中，热稳定性最好的是
    \begin{tasks}(4)
        \task $\ce{NaN3}$
        \task $\ce{AgN3}$
        \task $\ce{Pb(N3)2}$
        \task $\ce{Ba(N3)2}$
    \end{tasks}

    \item 下列说法正确的是
    \begin{tasks}(1)
        \task 碱性：$\ce{NH3}>\ce{N2H4}>\ce{NH2OH}$
        \task 还原性：$\ce{NH3}>\ce{N2H4}>\ce{NH2OH}$
        \task 熔点：$\ce{NH3}>\ce{N2H4}>\ce{NH2OH}$
        \task 热稳定性：$\ce{NH2OH}>\ce{N2H4}>\ce{NH3}$
    \end{tasks}

    \item 能和$\ce{FeSO4}$和$\ce{H2SO4}$发生棕色环反应的是
    \begin{tasks}(4)
        \task $\ce{NaNH2}$
        \task $\ce{NaN3}$
        \task $\ce{NaNO3}$
        \task $\ce{NaNO2}$
    \end{tasks}

    \item 某未知液中含有$\ce{K2SO4}$或$\ce{K2SO3}$，下列试剂中能鉴别他们的是
    \begin{tasks}(4)
        \task $\ce{NaCl(aq)}$
        \task $\ce{Br2(aq)}$
        \task $\ce{H2O2}$
        \task $\ce{KOH(aq)}$
    \end{tasks}

    \item 下列分子或离子中，中心原子杂化方式与$\ce{IF3}$相同的是
    \begin{tasks}(4)
        \task $\ce{IO3-}$
        \task $\ce{I3-}$
        \task $\ce{ICl2-}$
        \task $\ce{IF5}$
    \end{tasks}

    \item 在无氧条件下，用锌粉与亚硫酸氢钠和亚硫酸的混合溶液反应，得到
    \begin{tasks}(4)
        \task $\ce{Na2S}$
        \task $\ce{Na2S2O3}$
        \task $\ce{Na2S2O4}$
        \task $\ce{Na2S2O5}$
    \end{tasks}

    \item 向$\ce{AgNO3}$溶液中通入以下气体，没有$\ce{Ag}$单质沉淀生成的是
    \begin{tasks}(4)
        \task $\ce{NH3}$
        \task $\ce{N2H4}$
        \task $\ce{PH3}$
        \task $\ce{AsH3}$
    \end{tasks}

    \item 下列方程式书写正确的是(缺选项A、C)
    \begin{tasks}(1)
        \task （缺）
        \task $\ce{P4}$和氢氧化钠
        \task （缺）
        \task 将甲硅烷通入高锰酸钾
    \end{tasks}

    \item 下列配合物或配离子中，没有反馈$\pi$键的是
    \begin{tasks}(4)
        \task $[\ce{Cd(CN)6}]^{4-}$
        \task $[\ce{FeF6}]^{3-}$
        \task $[\ce{Pt(C2H4)Cl3}]^{-}$
        \task $[\ce{Fe(phen)3}]^{2+}$
    \end{tasks}
\end{enumerate}

\vspace{0.5cm}
\noindent\textbf{二、简要回答下列问题。本题共7小题，第11题15分，12~17题每题5分，共45分。}

\begin{enumerate}[label=\arabic*., itemsep=1.5ex]
    \item 解释或给出实验现象并写出必要的相关化学反应方程式：
    \begin{enumerate}[label=(\arabic*), noitemsep, topsep=0pt, partopsep=0pt]
        \item 向碱性$\ce{H2O2}$溶液中滴加$\ce{K2Cr2O7}$溶液得到黄色溶液；向酸性$\ce{H2O2}$溶液中滴加$\ce{K2Cr2O7}$溶液得到蓝色溶液，进而蓝色溶液逐渐变绿。
        \item 向氯化汞溶液中滴加$\ce{SnCl2}$溶液，先生成白色沉淀，随着$\ce{SnCl2}$溶液滴加至过量，沉淀逐渐变灰、变黑；滤去沉淀，再通入$\ce{H2S}$，得到黄色沉淀。
        \item $\ce{CoCl2}$溶液与过量$\ce{NaOH}$溶液生成粉红色沉淀，再用$\ce{NaClO}$溶液处理，粉红色沉淀变为棕黑色。该棕黑色沉淀溶于浓盐酸，所得溶液加水稀释后溶液变为粉红色。
    \end{enumerate}

    \item 比较$\ce{Na2CO3}$、$\ce{MgCO3}$、$\ce{K2CO3}$、$\ce{MnCO3}$和$\ce{PbCO3}$的热稳定性，并简述理由。

    \item 以海水、$\ce{AgNO3}$、$\ce{Fe}$和$\ce{Cl2}$为原料制备单质碘，写出各主要步骤的化学反应方程式。

    \item 给出并解释$\ce{SOF2}$、$\ce{SOCl2}$和$\ce{SOBr2}$的空间构型与$\ce{S-O}$键长递变规律。

    \item 有$\ce{NH4Cl}$、$\ce{Cd(NO3)2}$、$\ce{AgNO3}$、$\ce{ZnSO4}$、$\ce{Hg(NO3)2}$五瓶无色溶液失落标签，请选用一种试剂，将它们进行区别，分别写出所对应现象和反应方程式。

    \item 用晶体场理论分析比较$[\ce{Co(H2O)6}]^{2+}$、$[\ce{Co(NH3)6}]^{2+}$、$[\ce{Co(CN)6}]^{4-}$的还原性强弱。

    \item 硫代硫酸钠在药剂中常用作解毒剂，可解卤素单质、重金属离子及氰化物中毒。试说明其原理，并写出反应方程式。
\end{enumerate}

\vspace{0.5cm}
\noindent\textbf{三、推断题：本题共2小题，每小题5分，共10分。}

\begin{enumerate}[label=\arabic*., itemsep=1.5ex]
    \item $\ce{A}$不稳定易分解，分解得到$\ce{B}$和$\ce{NaCl}$，$\ce{B}$的水溶液电解可得$\ce{H2}$和$\ce{C}$。$\ce{B}$中加入浓盐酸得到黄色气态化合物$\ce{D}$，$\ce{D}$是具有漂白、消毒作用的强氧化剂，且显顺磁性。$\ce{D}$和$\ce{Na2O2}$反应生成$\ce{A}$和助燃性气体$\ce{E}$。

    \item 白色固体$\ce{A}$加热分解得到黄色固体$\ce{B}$和红棕色混合气体，将混合气体通入冰盐水中冷却得到无色液体$\ce{C}$以及一种气体单质。$\ce{A}$溶于稀硝酸后，若加入$\ce{NaOH}$溶液则得到白色沉淀$\ce{D}$；若加入$\ce{KI}$溶液则得到金黄色沉淀$\ce{E}$，$\ce{E}$可溶于热水或过量的$\ce{KI}$溶液。
\end{enumerate}

\vspace{1cm}
\begin{center}
    {\Large \textbf{物理化学部分(共75分)}}
\end{center}

\vspace{0.5cm}
\noindent\textbf{四、选择题：本题共 12 小题，每小题 2 分，共 24 分。在每小题给出的四个选项中，只有一项是符合题目要求的。}

\begin{enumerate}[label=\arabic*., itemsep=1.5ex]
    \item 已知 373.15 K 时，纯液体 A 的饱和蒸气压为 133.32 kPa，另一液体 B 可与 A 构成理想液体混合物。当 A 在溶液中的物质的量分数为 1/2 时，A 在气相中的物质的量分数为 2/3。则在 373.15 K 时，纯液体 B 的饱和蒸气压应为
    \begin{tasks}(4)
        \task 66.66 kPa
        \task 88.88 kPa
        \task 133.32 kPa
        \task 266.64 kPa
    \end{tasks}

    \item 在恒温恒压下形成理想液体混合物的混合吉布斯自由能 $\Delta_{\text{mix}}G \neq 0$，恒压下 $\Delta_{\text{mix}}G$ 对温度 $T$ 进行微商，则：
    \begin{tasks}(2)
        \task $(\partial \Delta_{\text{mix}}G / \partial T)_p < 0$
        \task $(\partial \Delta_{\text{mix}}G / \partial T)_p > 0$
        \task $(\partial \Delta_{\text{mix}}G / \partial T)_p = 0$
        \task $(\partial \Delta_{\text{mix}}G / \partial T)_p \neq 0$
    \end{tasks}

    \item 测定电池电动势时，标准电池的作用是
    \begin{tasks}(4)
        \task 提供标准电极电势
        \task 提供稳定电流
        \task 提供标准电位
        \task 提供稳定电压
    \end{tasks}

    \item 用对消法测定由电极 \ce{Ag(s)} $\mid$ \ce{AgNO3(aq)} 与电极 \ce{Ag(s), AgCl(s)} $\mid$ \ce{KCl(aq)} 组成的电池的电动势，下列哪一项是不能采用的：
    \begin{tasks}(4)
        \task 饱和 \ce{KCl} 盐桥
        \task 标准电池
        \task 电位计
        \task 直流检偏计
    \end{tasks}

    \item \textcolor{red}{(缺)}\\
    密闭容器，一边是 $101.325\ \text{kPa}$，$2\ \text{dm}^3$，$373.15\ \text{K}$，另一边为 $101.325\ \text{kPa}$，$1\ \text{dm}^3$，$373.15\ \text{K}$，现在将中间隔板撤掉，问下列哪个是正确的
    \begin{tasks}(4)
        \task $\Delta U$ \dots
        \task \dots
        \task \dots
        \task \dots
    \end{tasks}

    \item 在 298 K、$p^\ominus$ 下，两瓶含萘的苯溶液，第一瓶为 $2\ \text{dm}^3$（溶有 $0.5\ \text{mol}$ 萘），第二瓶为 $1\ \text{dm}^3$（溶有 $0.25\ \text{mol}$ 萘），若以 $\mu_1$ 和 $\mu_2$ 分别表示两瓶中萘的化学势，则
    \begin{tasks}(4)
        \task $\mu_1 = 10 \mu_2$ 
        \task $\mu_1 = 2 \mu_2$ 
        \task $\mu_1 = 1/2 \mu_2$ 
        \task $\mu_1 = \mu_2$
    \end{tasks}
    \item \textcolor{red}{(缺，但似乎不影响做题)}\\
    在密闭容器中，加入草酸钙，通入氮气排除空气后，加热草酸钙分解问 $C$ 和 $f$ 等于多少
    \begin{tasks}(4)
        \task $C=2, f=1$
        \task $C=2, f=2$
        \task $C=3, f=1$
        \task $C=3, f=2$
    \end{tasks}

    \item \textcolor{red}{(缺)}\\
    $^{14}\text{C}$ 衰变反应，他的 $t_{1/2}$ 和浓度的。。。半衰期公式（$C_0$）

    \item 某实际气体反应在温度为 500 K、压力为 $202.6\times 10^2\ \text{kPa}$ 下的平衡常数 $K_f=2$，则该反应在 500 K、$20.26\ \text{kPa}$ 下反应的平衡常数 $K_p$ 为
    \begin{tasks}(4)
        \task 2
        \task $>2$
        \task $<2$
        \task $\ge 2$
    \end{tasks}

    \item \textcolor{red}{(缺)}

    \item \textcolor{red}{(缺)}

    \item 在希托夫（Hittorf）法测迁移数的实验中，用 Ag 电极电解 \ce{AgNO3} 溶液，测出在阳极部 \ce{AgNO3} 增加了 $x\ (\text{mol})$，而串联在电路中的 Ag 库仑计上有 $y\ \text{mol}$ 的 Ag 析出，则 \ce{Ag+} 迁移数为（ \quad ）。
    \begin{tasks}(4)
        \task $\frac{x}{y}$
        \task $\frac{y}{x}$
        \task $\frac{x-y}{x}$
        \task $\frac{y-x}{y}$
    \end{tasks}
\end{enumerate}

\vspace{0.5cm}
\noindent\textbf{五、填空题：本题共 10 小题，每小题 1 分，共 10 分。}

\begin{enumerate}[label=\arabic*., itemsep=1.5ex]
    \item \textcolor{red}{(缺)} $5\ \text{mol}\ \ce{CO2}$，从 300K 加热到 350K，问其焓变值\underline{\qquad\qquad\qquad}（填增大、减小或不变）。
    
    \item 单质碘的熔点随压强的\underline{\qquad\qquad\qquad}（填增大、减小或不变）而降低。
    
    \item 临界状态下，水的气液两相的密度相对大小为\underline{\qquad\qquad\qquad}。
    
    \item \textcolor{red}{(缺)}
    
    \item 完全互溶双液系由 A 和 B 两组分组成，在 $x_A=0.4$ 存在最高恒沸混合物。现在想要通过精馏塔塔顶馏出 B，则混合物的浓度范围是\underline{\qquad\qquad\qquad}。
    
    \item \ce{Cl2}、\ce{H2} 在光照下反应一段时间后爆炸，这属于链反应中的\underline{\qquad\qquad\qquad}。
    
    \item 0.01 和 0.1 电导率 \textcolor{red}{(缺)}
    
    \item \textcolor{red}{(缺)}
    
    \item \textcolor{red}{(缺)} 氢超电势？
    
    \item 温度升高时，活化分子的比例\underline{\qquad\qquad\qquad}（填增大、减小或不变）。
\end{enumerate}

\vspace{0.5cm}
\noindent\textbf{六、简答题：简要回答下列问题。本题共 3 小题，每小题 5 分，共 15 分。}

\begin{enumerate}[label=\arabic*., itemsep=1.5ex]
    \item 用电化学方法算 \ce{AgI} 的活度积常数 $K_{sp}$
    
    \item 用物理化学的方法，说明旱和涝的土壤对植物生长的影响
    
    \item 下图所示是一个平行反应，其中 C 是所需要的产物，而 D 是副产物。假定两个反应的指前因子近似相同且不随温度变化而变化，能通过升温提高 C 的产率吗？试以 $\ln k \sim 1/T$ 对两反应作关联图说明。
    \[
    \ce{A + B} 
    \left\{
    \begin{array}{l}
        \xrightarrow{k_1,\ E_{a,1}} \ce{C} \\[1ex]
        \xrightarrow{k_2,\ E_{a,2}} \ce{D}
    \end{array}
    \right.
    \]
\end{enumerate}

\vspace{0.5cm}
\noindent\textbf{七、解答题：回答下列问题，解答应给出详细步骤。本题共 3 小题，第 1 题 8 分，第 2 题 8 分，第 3 题 10 分，共 26 分}

\begin{enumerate}[label=\arabic*., itemsep=1.5ex]
    \item 若 $1.0\ \text{mol}$ 正交硫转变为单斜硫时，体积增加 $3.53 \times 10^{-6}\ \text{m}^3$ (设此值不随温度、压而改变)，试问在 $100^{\circ}\text{C}$、$100p^\ominus$ 下，硫的那一种晶型较稳定？\\
    已知：在 298K、$101325\ \text{Pa}$ 下，正交硫和单斜硫的摩尔燃烧焓分别为 $-296.7\ \text{kJ}\cdot\text{mol}^{-1}$ 和 $297.1\ \text{kJ}\cdot\text{mol}^{-1}$。在 $p^\ominus$ 下，两种晶型的正常转化温度为 $96.7^{\circ}\text{C}$，假定硫的两种晶型的 $C_{p,m}$ 相等。
    
    \item 有人提出 $\ce{H2(g) + I2(g) -> 2HI(g)}$ 的反应机理如下：
    \begin{align*}
        \ce{I2} &\xrightarrow{k_1} \ce{2I} \\
        \ce{2I} &\xrightarrow{k_2} \ce{I2} \\
        \ce{2I + H2} &\xrightarrow{k_3} \ce{2HI}
    \end{align*}
    （1）假定 I 处于稳态，试导出 \ce{I2} 消失的速率方程 $-\frac{\mathrm{d}[\ce{I2}]}{\mathrm{d}t}$ \\
    （2）基元反应活化能和总反应表观活化能关系推导
    
    \item 用 Pt 做电极电解 \ce{SnCl2} 水溶液，在阴极上因 \ce{H2} 有超电势故只析出 \ce{Sn(s)}，在阳极上析出 \ce{O2}。已知 $a(\ce{Sn^{2+}})=0.10$，$a(\ce{H+})=0.010$；氧在阳极上析出的超电势为 0.500 V。\\
    已知：$\varphi^\ominus(\ce{Sn^{2+}/Sn}) = -0.140\ \text{V}$，$\varphi^\ominus(\ce{O2/H2O}) = 1.23\ \text{V}$。
    \begin{enumerate}[label=（\arabic*）]
        \item 写出电池反应和电极反应，并计算题设体系的实际分解电压；
        \item 氢在阴极上析出时的超电势为 0.500 V，计算要使 $a(\ce{Sn^{2+}})$ 降至何值时，才开始析出氢气？
    \end{enumerate}
\end{enumerate}